\begin{itemize}
\item Hourly concentrations spiked above 100 \(\mu\)g/m\(^3\) during 6.5\% of monitored hours.
\item Low-income (income\_quart = bin 1) households experience 114.2\% higher indoor PM2.5 than high-income households (income\_quart = bin 4).
\item We estimate that households in areas with more than two waste-burning events per week (10.7\% of the sample) experienced both higher average indoor PM2.5 concentrations and a higher probability of spikes in indoor PM2.5.
\item The probability of measuring spikes in indoor pollution levels was also \textcolor{magenta}{341.8\%} higher in smoking households than in non-smoking households.
\item 57.7\% of our sample lives within 500 meters from a primary road (dist\_primary).
\item Overall, we estimate that differences in infiltration explain 14.4\% of the difference in average indoor PM2.5 between the lowest and highest income quartile.
\item Overall, we estimate that smoking differences explain 47.6\% of the difference in average indoor PM2.5 between the lowest and highest income quartile.
\item We estimate that differences in other hyperlocal pollution sources explain the remaining 38.0\% of the difference in average indoor PM2.5 between the lowest and highest income quartile.
\end{itemize}
